\documentclass[11pt,a4paper,oneside]{article}
\usepackage[utf8]{inputenc}
\usepackage[spanish]{babel}
\usepackage{amsmath}
\usepackage{amsfonts}
\usepackage{amssymb}
\usepackage{graphicx}
\usepackage[colorlinks=true, linkcolor=blue, urlcolor=cyan]{hyperref}
\usepackage{listings}
\usepackage{xcolor}
\usepackage{geometry}
\geometry{margin=2.5cm}

\definecolor{codegreen}{rgb}{0,0.6,0}
\definecolor{codegray}{rgb}{0.5,0.5,0.5}
\definecolor{codepurple}{rgb}{0.58,0,0.82}
\definecolor{backcolour}{rgb}{0.95,0.95,0.92}

\lstdefinestyle{mystyle}{
    backgroundcolor=\color{backcolour},   
    commentstyle=\color{codegreen},
    keywordstyle=\color{magenta},
    numberstyle=\tiny\color{codegray},
    stringstyle=\color{codepurple},
    basicstyle=\ttfamily\footnotesize,
    breakatwhitespace=false,         
    breaklines=true,                 
    captionpos=b,                    
    keepspaces=true,                 
    numbers=left,                    
    numbersep=5pt,                  
    showspaces=false,                
    showstringspaces=false,
    showtabs=false,                  
    tabsize=2
}

\lstset{style=mystyle}

\title{MOOVIA | Reporte de Arquitectura Técnica y Algoritmos del Sensor}
\author{Ingeniería de Software - MOOVIA App}
\date{\today}

\begin{document}

\maketitle
\tableofcontents
\newpage

\section{Introducción}
Este documento describe la arquitectura técnica de la plataforma MOOVIA para el seguimiento de levantamientos en halterofilia. El sistema utiliza un sensor IMU (Inertial Measurement Unit) ICM-42688-P controlado por un microcontrolador WBZ351 para capturar y procesar datos biocinemáticos en tiempo real.

\section{Arquitectura del Software (Monorepo)}
La solución está estructurada como un monorepo utilizando \textbf{npm workspaces}, lo que permite separar la lógica de negocio de la interfaz de usuario y compartir algoritmos entre diferentes plataformas (Mobile y Web).

\begin{itemize}
    \item \textbf{apps/mobile}: Aplicación React Native (Expo) que captura datos vía Bluetooth Low Energy (BLE).
    \item \textbf{apps/web-test}: Herramienta de visualización técnica (Vite + TS) para validación de algoritmos offline.
    \item \textbf{packages/sensor-core}: Paquete de TypeScript puro que contiene la lógica de procesamiento (Decoders, Madgwick, Kalman).
\end{itemize}

\section{Pipeline de Procesamiento de Datos}
El flujo de datos sigue un proceso de múltiples etapas desde la captura hasta la obtención de métricas finales:

\subsection{1. Captura y Decodificación (BLE)}
El sensor envía paquetes de 211 bytes vía notificaciones BLE. Cada paquete contiene 15 muestras.
\begin{itemize}
    \item \textbf{Frecuencia (ODR)}: 1000 Hz (1 ms entre muestras).
    \item \textbf{Formato}: Datos en \textit{little-endian int16}.
    \item \textbf{Conversión Física}: Los valores crudos se convierten a unidades internacionales usando las resoluciones:
    \begin{equation}
        a_{physical} = \frac{a_{raw}}{4096} \, g, \quad w_{physical} = \frac{w_{raw}}{32.8} \, dps
    \end{equation}
\end{itemize}

\subsection{2. Estimación de Orientación (Filtro Madgwick)}
Para eliminar el componente de gravedad del acelerómetro y conocer la inclinación de la barra, se utiliza un filtro de orientación de Madgwick. Este filtro fusiona datos del giróscopo y acelerómetro para obtener un cuaternión de orientación estable, minimizando el drift.

\subsection{3. Estimación de Cinemática (Filtro Kalman)}
Se utiliza un filtro de Kalman de 3 estados por eje para estimar:
\begin{equation}
    x = [p, v, b]^T
\end{equation}
Donde $p$ es la posición, $v$ la velocidad y $b$ el bias del acelerómetro acumulado. El filtro aplica \textbf{ZUPT (Zero Velocity Update)} cuando detecta que el sensor está estacionario, corrigiendo el drift de integración.

\section{Algoritmos de Métricas Propias}
El \texttt{TrajectoryService.ts} calcula métricas clave para el entrenamiento de fuerza:

\subsection{Velocidad Media Propulsiva (VMP)}
Basado en los estudios de Sánchez-Medina y González-Badillo, la VMP es la media de la velocidad durante la fase concéntrica (subida) del levantamiento, excluyendo la fase de frenado donde la aceleración es menor que la gravedad.
\begin{lstlisting}[language=Java]
// Pseudocodigo del calculo de VMP
if (vZ > 0.05 && aZ > -9.81) {
    sumVelocity += vZ;
    count++;
}
\end{lstlisting}

\subsection{Aceleración Lineal Pico}
Es el valor máximo de la magnitud del vector de aceleración lineal (sin gravedad) durante el movimiento.
\begin{equation}
    a_{peak} = \max \sqrt{a_x^2 + a_y^2 + a_z^2}
\end{equation}

\subsection{Altura Máxima}
Diferencia máxima en la coordenada Z con respecto al punto de inicio del levantamiento, calculada tras re-procesar la trayectoria para eliminar bias.

\section{Herramienta de Diagnóstico (Web Test Lab)}
Permite a los ingenieros cargar archivos JSON grabados por la aplicación móvil. La herramienta recrea el procesamiento completo paso a paso, permitiendo visualizar:
\begin{itemize}
    \item Graficos de aceleración neta.
    \item Integración de velocidad vertical.
    \item Trayectoria 3D de la barra.
\end{itemize}

\section{Conclusión}
La arquitectura desacoplada garantiza que las mejoras en el algoritmo de procesamiento en \texttt{sensor-core} se reflejen inmediatamente tanto en la herramienta de validación como en la aplicación móvil de producción.

\end{document}
